\section{Introduction}\label{sec:introduction}
% Budget: 1.5 pages.
\todo{Motivation: AI-based perception in automated driving; monitors as architectural safety mechanism; who guards the guards; where to place monitors.}

\todo{Problem statement and gap.}

\paragraph{Research Questions.}
\begin{description}
  \item[RQ1] What share of realistic safety requirements on a perception pipeline can be expressed in the monitorable contract fragment?
  \item[RQ2] Do component-level monitors detect violations of the system-level contract, and with what delay and false-alarm rate?
  \item[RQ3] Do separate assumption and guarantee monitors attribute a violation to the correct component or to its environment?
\end{description}

\paragraph{Contributions.} We provide the following contributions:
\begin{enumerate}
  \item[(i)] a contract framework over timed streams with refinement, composition, and a prefix semantics that fixes what can be decided at runtime;
  \item[(ii)] a compositional monitor synthesis along the refinement hierarchy, with mechanised proofs of monitor correctness, of the equivalence of centralised and distributed monitoring, and of the preservation of system-level guarantees under refinement;
  \item[(iii)] contract patterns for AI-based perception components;
  \item[(iv)] an evaluation on a perception pipeline using fault injection.
\end{enumerate}
